\documentclass[12pt,a4paper]{article}
\usepackage[margin=2.5cm]{geometry}
\begin{document}

\begin{flushright}
Zhilei Chen\\
Guangdong Peizheng College\\
53 Peizheng Rd., Chini Town, Huadu\\
Guangzhou, Guangdong 510830, China\\
2604513@peizheng.edu.cn\\
\today
\end{flushright}

\vspace{1cm}

\noindent
The Editor\\
\textit{International Journal of Fatigue}\\
Elsevier

\vspace{0.5cm}

\noindent Dear Editor,

\medskip

\noindent I am submitting the manuscript entitled ``Method-Induced
Uncertainty in Gear Bending Fatigue: Factorial Decomposition of
Size-and-Surface Factor Sensitivity Under a Common ISO Stress
Baseline'' for consideration for publication in the
\textit{International Journal of Fatigue}.

This paper addresses a practical problem in gear design: when three
qualified engineers compute the bending fatigue life of the same spur
gear---identical geometry, material, and load---their predictions can
span five orders of magnitude. The divergence originates not from
material variability but from the methodological choices each engineer
legitimately makes under the letter of published standards.

The manuscript makes three contributions:

\begin{enumerate}
  \item \textbf{First full-factorial variance decomposition across five
    simultaneous engineering choices in gear fatigue.} No prior study
    has ranked the contributions of standard selection, mean-stress
    correction method, surface roughness grade, S--N curve source, and
    cumulative damage rule within a single controlled experimental
    design. The analysis employs a censored-likelihood AFT framework
    that preserves information from the 28 run-out observations
    (19\% of the 144-combination dataset)---unlike conventional ANOVA
    or Sobol' sensitivity analysis, which must discard them.
  \item \textbf{Demonstration that method-induced uncertainty dominates
    the prediction budget.} A quantitative benchmark against published
    42CrMo4 gear fatigue test data (Wang et al., 2023) shows that the
    method-induced $\log_{10}(N_f)$ spread (2.1--3.3~dex) is 2--3
    times larger than the intrinsic material scatter ($\sim$1~dex) at
    a fixed stress level. The three dominant factors (mean-stress
    correction, size-and-surface factor formulation, surface roughness)
    together account for 88.6\% of total log-variance.
  \item \textbf{Translation of numerical results into actionable
    engineering guidance.} The paper provides cross-standard empirical
    correction formulas, a design risk map ranking methodological
    choices by conservatism direction, consequence-classified
    uncertainty thresholds, and a step-by-step design protocol for
    controlling prediction scatter.
\end{enumerate}

A complete reproducibility package is submitted as Supplementary
Material, including all Python source code, the 144-combination sweep
data, and a one-click reproduction script (\texttt{bash reproduce.sh}).

This manuscript has not been published previously and is not under
consideration elsewhere. All authors have approved the submission. The
author declares no competing interests. Generative AI tools
(Claude/Anthropic) were used for language editing and code
development; all scientific content, analysis decisions, and numerical
results are the author's own.

I believe this work fits the scope of IJF as a computational
uncertainty quantification study that addresses a recognized practical
problem in gear fatigue assessment. I look forward to the reviewers'
evaluation.

\vspace{0.5cm}

\noindent Sincerely,

\vspace{0.5cm}

\noindent Zhilei Chen

\end{document}
